\documentclass[CJK]{beamer}
\usepackage{CJKutf8}
\usepackage{beamerthemesplit}
\usetheme{Malmoe}
\useoutertheme[footline=authortitle]{miniframes}
\usepackage{amsmath}
\usepackage{amssymb}
\usepackage{graphicx}
\usepackage{color}
\usepackage{slashed}
\usepackage{simplewick}
\graphicspath{{../figures/}}
\def\be{\begin{equation}}
\def\ee{\nonumber\end{equation}}
\def\bea{\begin{eqnarray}}
\def\eea{\nonumber\end{eqnarray}}
\def\ii{{\dot{\imath}}}
\def\bch{\begin{CJK}{UTF8}{gbsn}}
\def\ech{\end{CJK}}
\def\bex{\begin{minipage}{0.3\textwidth}\includegraphics[width=1in]{jugelizi.png}\end{minipage}\begin{minipage}{0.6\textwidth}}
\def\eex{\end{minipage}}
\def\chtitle#1{\frametitle{\bch#1\ech}}
\def\skipline{{\vskip0.1in}}
\def\skiplines{{\vskip0.2in}}
\def\lagr{{\mathcal{L}}}
\def\hamil{{\mathcal{H}}}
\def\vecv{{\mathbf{v}}}
\def\vecx{{\mathbf{x}}}
\def\veck{{\mathbf{k}}}
\def\vecp{{\mathbf{p}}}
\def\vecn{{\mathbf{n}}}
\def\vecA{{\mathbf{A}}}
\def\vecP{{\mathbf{P}}}
\def\vecsigma{{\mathbf{\sigma}}}
\def\hatJn{{\hat{J_\vecn}}}
\def\hatJx{{\hat{J_x}}}
\def\hatJy{{\hat{J_y}}}
\def\hatJz{{\hat{J_z}}}
\def\hatj#1{\hat{J_{#1}}}
\def\hatphi{{\hat{\phi}}}
\def\hatq{{\hat{q}}}
\def\hatpi{{\hat{\pi}}}
\def\vel{\upsilon}
\def\Dint{{\mathcal{D}}}
\def\adag{{\hat{a}^\dagger}}
\def\bdag{{\hat{b}^\dagger}}
\def\cdag{{\hat{c}^\dagger}}
\def\ddag{{\hat{d}^\dagger}}
\def\hata{{\hat{a}}}
\def\hatb{{\hat{b}}}
\def\hatc{{\hat{c}}}
\def\hatd{{\hat{d}}}
\def\hatN{{\hat{N}}}
\def\hatH{{\hat{H}}}
\def\hatp{{\hat{p}}}
\def\Fup{{F^{\mu\nu}}}
\def\Fdown{{F_{\mu\nu}}}
\def\newl{\nonumber \\}
\def\SIkm{\mathrm{km}}
\def\SIyr{\mathrm{yr}}
\def\SIGyr{\mathrm{Gyr}}
\def\SIeV{\mathrm{eV}}
\def\SIGeV{\mathrm{GeV}}
\def\SIm{\mathrm{m}}
\def\SIcm{\mathrm{cm}}
\def\SIJ{\mathrm{J}}
\def\SIs{\mathrm{s}}
\def\SIkg{\mathrm{kg}}
\def\SIg{\mathrm{g}}
\def\vece{\mathrm{e}}
\def\bmat#1{\left(\begin{array}{#1}}
\def\emat{\end{array}\right)}
\def\bcase#1{\left\{\begin{array}{#1}}
\def\ecase{\end{array}\right.}
\def\calM{{\mathcal{M}}}
\def\calT{{\mathcal{T}}}
\def\calR{{\mathcal{R}}}
\def\barpsi{\bar{\psi}}
\def\baru{\bar{u}}
\def\barv{\bar{\upsilon}}
\def\bmini#1{\begin{minipage}{#1\textwidth}}
\def\emini{\end{minipage}}
\def\qeq{\stackrel{?}{=}}
\def\torder#1{\mathcal{T}\left(#1\right)}
\def\rorder#1{\mathcal{R}\left(#1\right)}


\title{Quantum Field Theory I \\ Homework 6 solution}
  \author{}
  \date{}


\begin{document}

\begin{frame}
 
\begin{center}
\begin{Large}
\bch
量子场论 I 

{\vskip 0.3in}

第六次课后作业参考答案
\skipline
\skipline

如发现参考答案有错误请不吝告知（微信zhiqihuang或邮箱huangzhq25@sysu.edu.cn)
\ech
\end{Large}
\end{center}

\vskip 0.2in

\bch
课件下载
\ech
https://github.com/zqhuang/SYSU\_QFTI

\end{frame}


\begin{frame}
\chtitle{第1题 题目和思路}
\bch
{\small
证明下述矩阵的迹为零：
\begin{itemize}
\item{$\slashed{a}\slashed{b}\gamma^0$}
\item{$\slashed{a}\slashed{b}\gamma^5$}
\item{$\slashed{a}\slashed{b}\slashed{c}$}
\item{$\slashed{a}\slashed{b}\slashed{c}\gamma^5$}
\item{$\slashed{a}\slashed{b}\slashed{c}\slashed{d}-\slashed{d}\slashed{c}\slashed{b}\slashed{a}$}
\end{itemize}

\skipline
思路：用$\gamma$矩阵的零迹定理和Feynman符号迹的倒排定理
}


\ech
\end{frame}

\begin{frame}
\chtitle{第1题解答}
\bch
{\small
\begin{itemize}
\item{$\slashed{a}\slashed{b}\gamma^0$展开式的每一项都是$\gamma^\mu\gamma^\nu\gamma^0$的形式($0\le \mu,\nu\le 3$)。根据零迹定理（排除$\gamma^5$），$\gamma^\mu\gamma^\nu\gamma^0$的迹为零。}
\item{$\slashed{a}\slashed{b}\gamma^5$的展开式每一项都是$\gamma^\mu\gamma^\nu\gamma^5$的形式($0\le \mu,\nu\le 3$)。 根据零迹定理（取某个$\alpha \ne \mu,\nu, 5$, 排除$\gamma^\alpha$） $\gamma^\mu\gamma^\nu\gamma^5$的迹为零}
\item{$\slashed{a}\slashed{b}\slashed{c}$的展开式每一项都是$\gamma^\mu\gamma^\nu\gamma^\lambda$的形式($0\le \mu,\nu,\lambda\le 3$)。根据零迹定理（排除$\gamma^5$)，$\gamma^\mu\gamma^\nu\gamma^\lambda$的迹为零。}
\item{$\slashed{a}\slashed{b}\slashed{c}\gamma^5$展开式的每一项是$\gamma^\mu\gamma^\nu\gamma^\lambda\gamma^5$的形式($0\le \mu,\nu, \lambda \le 3$)。根据零迹定理（排除$\gamma^5$)，$\gamma^\mu\gamma^\nu\gamma^\lambda\gamma^5$的迹为零。}
\item{根据Feynman符号的迹的倒排定理，$\slashed{a}\slashed{b}\slashed{c}\slashed{d}-\slashed{d}\slashed{c}\slashed{b}\slashed{a}$的迹为零。}
\end{itemize}

}


\ech
\end{frame}

\begin{frame}
\chtitle{第2题 题目和思路}
\bch
证明下述恒等式：
\begin{itemize}
\item{$\slashed{a}\slashed{b} + \slashed{b}\slashed{a} = 2 ab$}
\item{$\trof{\slashed{a}\gamma^\mu} = 4 a^\mu$}
\end{itemize}

\skipline
思路:利用Feynman符号和$\gamma$矩阵的定义的定义即可证明
\ech
\end{frame}

\begin{frame}
\chtitle{第2题解答}
\bch
\begin{itemize}
\item{$\slashed{a}\slashed{b} + \slashed{b}\slashed{a} = a_\mu b_\nu (\gamma^\mu\gamma^\nu+\gamma^\nu\gamma^\mu) = 2 g^{\mu\nu}a_\mu b_\nu = 2 ab$}
\item{利用迹可以轮换乘积次序的性质，$\trof{\gamma^\nu\gamma^\mu} =\trof{\gamma^\mu\gamma^\nu} =\frac{1}{2}\trof{\gamma^\mu\gamma^\nu + \gamma^\nu\gamma^\mu} = g^{\mu\nu}\trof{I_{4\times 4}} = 4 g^{\mu\nu}$

所以 $\trof{\slashed{a}\gamma^\mu} = a_\nu \trof{\gamma^\nu\gamma^\mu} =  4 a_\nu g^{\mu\nu} = 4 a^\mu$}
\end{itemize}

\ech
\end{frame}

\begin{frame}
\chtitle{第3题 题目和思路}
\bch 
把
$$\trof{\slashed{a}\gamma^\mu\slashed{b}\slashed{c}\slashed{d}\gamma_\mu}$$
化简为只含矢量内积的最简形式。

\skipline
思路：利用$\gamma$矩阵形式下标矩阵的性质和Feynman符号迹的展开定理
\ech
\end{frame}

\begin{frame}
\chtitle{第3题 第一种解答}
\bch
利用$\gamma^\mu\gamma^\alpha\gamma^\beta\gamma^\lambda\gamma_\mu = -2 \gamma^\lambda\gamma^\beta\gamma^\alpha$即可得到$\trof{\gamma^\mu\slashed{b}\slashed{c}\slashed{d}\gamma_\mu} = -2\slashed{d}\slashed{c}\slashed{b}$，所以
$$\trof{\slashed{a}\gamma^\mu\slashed{b}\slashed{c}\slashed{d}\gamma_\mu} = -2\trof{\slashed{a}\slashed{d}\slashed{c}\slashed{b}} = -8 (ad)(bc)   +8 (ac)(bd) -8 (ab)(cd) $$
\ech
\end{frame}


\begin{frame}
\chtitle{第3题 第二种解答}
\bch
利用$\gamma_\mu\gamma^\alpha \gamma^\mu = -2 \gamma^\alpha$即可得到$\trof{\gamma^\mu\slashed{a}\gamma_\mu} = -2\slashed{a}$，所以
\bea
\trof{\slashed{a}\gamma^\mu\slashed{b}\slashed{c}\slashed{d}\gamma_\mu} &=& \trof{\gamma_\mu\slashed{a}\gamma^\mu\slashed{b}\slashed{c}\slashed{d}} \newl
&=& -2 \trof{\slashed{a}\slashed{b}\slashed{c}\slashed{d}} \newl
&=& -8(ab)(cd)   + 8(ac)(bd) -8(ad)(bc)
\eea
\ech
\end{frame}


\begin{frame}
\chtitle{第4题 题目和思路}
\bch
设$p$为电子的四维动量，$k$为光子的四维动量，化简
$$\trof{ (\slashed{p}+m)\gamma^\mu\frac{1}{\slashed{p}+\slashed{k}-m}\gamma_\mu}$$

\skipline
思路：利用第19课所讲的技巧
 
\ech
\end{frame}

\begin{frame}
\chtitle{第4题 第一种解答}
\bch
{\scriptsize
\bea
 \trof{ (\slashed{p}+m)\gamma^\mu\frac{1}{\slashed{p}+\slashed{k}-m}\gamma_\mu} &= &\trof{ (\slashed{p}+m)\gamma^\mu\frac{\slashed{p}+\slashed{k}+m}{(p+k)^2-m}\gamma_\mu} \newl
&=& \frac{1}{2pk}\trof{ (\slashed{p}+m)\gamma^\mu(\slashed{p}+\slashed{k}+m)\gamma_\mu} 
\eea
利用$\trof{(\slp+m)\gamma^\mu(\slp+m) } = 2p^\mu(\slp+m)$,上式等于
\bea
&& \frac{1}{2pk}\trof{ (\slashed{p}+m)(2p^\mu\gamma_\mu + \gamma^\mu\slashed{k}\gamma_\mu)} \newl
&=& \frac{1}{2pk}\trof{ (\slashed{p}+m)(2\slashed{p} - 2\slashed{k})} \newl
&=& \frac{1}{pk}\trof{\slashed{p}(\slashed{p} - \slashed{k}) } \newl
&=& \frac{4m^2}{pk} - 4
\eea
在最后一步我们利用了$\trof{\slp^2} = 4 p^2 = 4 m^2$和$\trof{\slp\slk}= 4pk$
 }
\ech
\end{frame}


\begin{frame}
\chtitle{第4题 第二种解答}
\bch
{\scriptsize
\be
 \trof{ (\slashed{p}+m)\gamma^\mu\frac{1}{\slashed{p}+\slashed{k}-m}\gamma_\mu} = \trof{\gamma_\mu (\slashed{p}+m)\gamma^\mu\frac{1}{\slashed{p}+\slashed{k}-m}} 
\ee
利用$\gamma_\mu \slp\gamma^\mu = -2\slp$和$\gamma_\mu\gamma^\mu = 4$，上式等于
\bea
\trof{ (-2\slashed{p}+4m)\frac{1}{\slp+\slk-m}} &=&\trof{ (-2\slashed{p}+4m)\frac{\slashed{p}+\slashed{k}+m}{(p+k)^2-m}} \newl
&=& \frac{1}{pk}\trof{ (-\slashed{p}+2m)(\slashed{p}+\slashed{k}+m)} 
\eea
展开上式，因为奇数个feynman符号的乘积的迹为零，只要保留偶次项。上式等于
\be
\frac{1}{pk}\trof{-\slashed{p}(\slashed{p}+\slashed{k}) + 2m^2} = \frac{1}{pk}\trof{m^2-\slashed{p}\slashed{k}) } = \frac{4m^2}{pk} - 4
\ee
 }
\ech
\end{frame}


\begin{frame}
\chtitle{第5题 题目和思路}
\bch
在Compton散射过程中，取电子初始状态为静止，入射光子沿$z$轴方向的参照系。限定入射光子四维动量为$k^\mu=(\omega, 0, 0, \omega)$；出射光子能量为$\omega'$，方向限定在某固定方向$\vecn$附近的立体角$d\Omega$内；$\vecn$与$z$轴的夹角为$\theta$。当$d\Omega$很小时，散射截面与$d\Omega$成正比，它们之间的比称为微分散射截面：$\frac{d\sigma}{d\Omega}$。试根据课上所求的散射概率计算该微分散射截面(写成$\omega$, $\omega'$,和$\theta$的函数)。

\skipline
思路：在限定范围内对末态进行积分
\ech
\end{frame}

\begin{frame}
\chtitle{第5题 解答}
\bch
{\scriptsize

在课上的结果
\be
|\calM|^2 = \frac{q^4}{8\omega_p\omega_p'\omega_k\omega_k'}\left[\frac{pk'}{pk} + \frac{pk}{pk'} + 2m^2\left(\frac{1}{pk}-\frac{1}{pk'}\right) + m^4\left(\frac{1}{pk}-\frac{1}{pk'}\right)^2\right]
\ee
中代入$pk = m \omega$, $pk' = m \omega'$，得到
$$|\calM|^2 = \frac{q^4}{8 m (m+\omega-\omega')\omega\omega'}\left[\frac{\omega'}{\omega} + \frac{\omega}{\omega'} + 2m\left(\frac{1}{\omega}-\frac{1}{\omega'}\right) + m^2\left(\frac{1}{\omega}-\frac{1}{\omega'}\right)^2\right]$$
由能量动量守恒可以得到
$$m\left(\frac{1}{\omega}-\frac{1}{\omega'}\right) = \cos\theta - 1$$
所以
$$|\calM|^2 = \frac{q^4}{8 m (m+\omega-\omega')\omega\omega'}\left[\frac{\omega'}{\omega} + \frac{\omega}{\omega'} - \sin^2\theta\right]$$
}

\ech
\end{frame}

\begin{frame}
\chtitle{第5题 解答（续）}
\bch
{\scriptsize
同第14课的推导方法，但不对末态$d\Omega$积分：
\be
\frac{d\sigma}{d\Omega} = \frac{1}{(2\pi)^2} \int \omega'^2 d\omega' \int d^3\vecp' |\calM|^2\delta(p+k-p'-k')
\ee
注意因末态粒子不同，不需要除以$2!$因子。入射速度$\upsilon = 1$所以也可以略去。

对$\vecp'$积分，得到
\be
\frac{d\sigma}{d\Omega} =\frac{1}{(2\pi)^2} \int \omega'^2 d\omega' |\calM|^2\delta(m+\omega-\omega'-\sqrt{m^2+(\omega^2+\omega'^2-2\omega\omega'\cos\theta)}) 
\ee
对$\omega'$进行积分，利用复合$\delta$函数的积分公式：
$$\int g(x) \delta(f(x)) dx = \sum_{x^*:f(x^*)=0} g(x^*)\frac{1}{|f'(x^*)|}$$
和

\be
\left\vert\frac{d\left(m+\omega-\omega'-\sqrt{m^2+(\omega^2+\omega'^2-2\omega\omega'\cos\theta)} \right)}{d\omega'} \right\vert  = 1 + \frac{\omega'-\omega\cos\theta}{\sqrt{m^2+(\omega^2+\omega'^2-2\omega\omega'\cos\theta)}}
\ee
}
\ech
\end{frame}

\begin{frame}
\chtitle{第5题 解答（续）}
\bch
{\scriptsize
以及能量守恒 $\sqrt{m^2+(\omega^2+\omega'^2-2\omega\omega'\cos\theta)} = m+\omega-\omega'$，我们得到：
\bea
\frac{d\sigma}{d\Omega} &=& \frac{1}{(2\pi)^2}  \omega'^2  |\calM|^2 \frac{1}{1+\frac{\omega'-\omega\cos\theta}{m+\omega-\omega'}} \newl
&=& \frac{1}{(2\pi)^2}  \omega'^2  |\calM|^2 \frac{m+\omega-\omega'}{m+\omega(1-\cos\theta)} \newl
&=& \frac{1}{(2\pi)^2}  \omega'^2  |\calM|^2 \frac{m+\omega-\omega'}{m+\omega m(\frac{1}{\omega'}-\frac{1}{\omega})} \newl
&=& \frac{1}{(2\pi)^2}  \omega'^2  |\calM|^2 \frac{\omega'(m+\omega-\omega')}{m\omega} \newl
&=& \frac{q^4}{32\pi^2}  \frac{\omega'^2}{\omega^2m^2}\left[\frac{\omega'}{\omega} + \frac{\omega}{\omega'} - \sin^2\theta\right]   \newl
\eea


}

\ech
\end{frame}

\end{document}
